% Chapter introduction
\label{sec:introduction}
% **********************************
\section{Introduction}
% **********************************

Dans la vie d'un étudiant et d'un professionnel, on doit rédiger des rapports. \`A l'université on considère
par défaut que l'étudiant doit apprendre par lui-même à faire celà en cherchant des exemples. Dans des universités
étrangères cela fait partie d'un enseignement. 

On peut s'apercevoir que au final, après 5 années d'études, peu savent préparer et rédiger un beau rapport ou un beau compte-rendu
d'un travail donné.

Il m'a semblé donc important de préparer un document assez complet, mais pas exhausitif rappelant les points essentiels,
distribuables à tous les niveaux de la licence au master. 

Ce document est rédigé sur la base de mon expérience dans l'écriture de cours,
 de livres, d'articles, de rapports techniques.
 Il prend aussi en compte  l'expérience acquise par des expertises très nombreuses de projets scientifiques ou d'articles pour des revues internationales.

Enfin, mon expérience a été renforcée  au cours de ma carrière par la lecture  de centaines de rapports
de stage ou de projets, et un grand nombre de thèses de doctorat.

Enfin la lecture d'ouvrage de référence sur l'écriture de rapports joints dans la bibliographie a permis de compléter l'ensemble.
 
Ce rapport en latex pourra servir d'exemple (templates) modifiable, mais les conseils donnés sont aisément transposables, 
parfois même plus facilement dans d'autres types de logiciels comme MS Word ou OO Writer.
Bien que  Latex présente des facilités  bien pratiques à l'écriture de rapport il n'est pas préféré au sein du monde industriel
qui favorise largement MS Word. Rappelons que Latex est un logiciel libre et 
donc privilégié à l'université et dans le monde de la recherche.
 
Enfin pour clore cette section il faut préciser qu'on ne perd par plus de temps à faire bien les choses que de mal les faire.
Même souvent c'est le contraire, car bien faire les choses amènent des questions utiles pour améliorer le contenu.


% ********************************
\section{Compétences à acquérir}
% *******************************

La vocation des Masters est de fournir aux étudiants des connaissances scientifiques et techniques et des compétences 
dans le but de favoriser au mieux une future insertion professionnelle.
 On peut y ajouter l'introduction des méthodologies de travail, des méthodes de pensées qui permettront d'aider à une
 formation tout au long de la vie, en autonomie ou pas. 

Il existe suivant les domaines scientifiques, d'ingénierie, de professions visées, 
 un grand nombre de compétences à acquérir tout au long des études dans l'enseignement supérieure, 
avec des niveaux  de difficultés progressifs.

Une des plus importantes  compétences transversales,  commune  consiste à 
 savoir rédiger correctement un rapport de qualité  qui est un document écrit en général archivé et donc transmissible.

Le rapport est une trace, voir un signature de l'auteur qu'il laisse derrière lui,
et qui met en avant ses compétences ou parfois son manque de compétence. 

Les exigences en termes de qualité et de rendu ont augmenté progressivement au cours des années,
 notamment à cause de l'arrivée d'outils informatiques très performants pour l'aide à l'écriture et à la publication.

Il convient donc se de former, et d'apprendre en autonomie à rédiger un document de qualité. 

Des recherches sur internet montrent que finalement on trouve très peu de documents qui indiquent clairement 
aux étudiants comment rédiger un compte rendu de TP ou un rapport de stage ou de projet. 
 

Un premier document a été rédigé  sur l'écriture d'un compte-rendu de TP.
Je continue  cette démarche en proposant ce document sur la rédaction d'un rapport de projet ou de stage.

Les documents ne sont pas parfaits et finalement présentent mon point de vue personnel.
Certaines parties peuvent notamment être plus tournées vers la science ou l'ingénierie, et ne conviendront pas naturellement 
à des rapports
dans d'autres domaines comme la santé,  les sciences humaines, les sciences économiques ou le droit.

J'espère  que ce document sera partagé, et que par des retours pourront  participer à son amélioration.


On peut conclure cette section en rappelant quelques compétences à acquérir pour rédiger un rapport de qualité, au moins dans la forme:

\begin{itemize}
\item être capable de maîtriser un ou plusieurs  logiciels de publication (MS Word, OO Write, Latex). Cela inclus:
	\begin{itemize}
	\item le choix ou la définition d'un style global
	\item l'ajout de tableaux ou de figures avec numérotation et titre
	\item la gestion de haut et bas de page
	\item la gestion d'une bibliographie avec référencement automatique, associée  à une  base de donnée bibliographique
	\item la gestion de la mise en page (figures et tableau à la bonne position par exemple)
	\end{itemize} 
\item être capable d'utiliser (voir de maîtriser) des outils de traitement d'images (GIMP par exemple) et des outils de dessin vectoriel (INKSCAPE par exemple).
       Des exemples sont données par la suite.
\item être capable de rédiger dans un français correct (ou un anglais correct) 
\item être capable d'organiser ses réflexions, ses analyses pour produire un document structuré 
\item être en mesure d'utiliser les fonctions graphiques de logiciels ou de langage (Python, ...) afin de produire des figures de qualité
\item \alert{A COMPLETER} 
\end{itemize}



%**********************
\section{Les rapports}
%**********************

%****************************************
\subsection{Différents types de rapport}
%***************************************

Pour un étudiant on peut considérer trois types de rapports:
\begin{itemize}
\item un compte-rendu de TP  qui est un petit rapport sur une activité très contrôlée et très formatée par un énoncé relativement précis qui sert de guide.
En général le compte-rendu ne dépasse pas les 20 pages, et très souvent les 10 pages.
\item un rapport de projet, scientifique ou technique. C'est un rapport plus libre qui doit mettre en avant un travail en autonomie d'un étudiant. 
Sa longueur est entre 10 et 30 pages. 
\item un rapport de stage. Ce dernier est une version améliorée d'un rapport de projet. La durée d'un stage est de plusieurs mois. Pour cette durée
 on estime qu'une longueur d'une quarantaine de pages, hors annexes, est amplement suffisante.
\end{itemize}

La taille d'un rapport suivant son type  est rappelée dans le tableau \ref{table:report}. 
La dernière ligne du tableau donne un élément de réflexion sur sa structure primaire, et donc la  possibilité d'utiliser
 des chapitres qui ne sont pas nécessaires si le document contient trop peu de pages.

\begin{table}[htb]
\setstretch{1.2}
\small
\centering
\begin{tabular}{p{2cm} p{4.5cm} p{4.5cm} p{4.5cm}}
 \Xhline{2pt}
type &   comte-rendu de TP  & rapport de projet &  rapport de stage \\ \Xhline{1pt}
pages &  5 - 20              & 10 - 30          &  30 - 40 \\
Chapitre & Non               & Peut-être        & Probable \\  \Xhline{2pt}
\end{tabular}
\caption{Longueur de différents rapports et existence de chapitre.}
\label{table:report}
\end{table}

\textbf{Un certain nombre d'indications fournies dans le document sur la rédaction d'un compte-rendu de Travail Pratique\footnote{travail de laboratoire (labwork en anglais)} restent valables et sont reprises dans ce document.}


%*******************************************************
\subsection{Rapport scientifique ou rapport technique}
%*******************************************************

La différence majeure entre un rapport scientifique et un rapport technique 
se situe dans le contenu.
 Naturellement cela parfois se répercute  aussi sur la forme. 
 Un rapport scientifique est beaucoup plus difficile à écrire et souvent plus long qu'un rapport technique.

On va prendre un exemple classique dans le milieu de la recherche afin d'illustrer les différences fondamentales.

%*********************************
\subsubsection{Rapport technique}
%*********************************

Une entreprise contacte une équipe de recherche  pour mener des essais en soufflerie. 
On doit effectuer voir concevoir le montage,  les essais et rédiger un rapport technique.
Ce rapport décrit le montage, explique comment les essais ont été menés et fournit les résultats des essais 
ainsi que des statistiques permettant d'évaluer les erreurs de mesure.

Un rapport technique contient donc uniquement des données techniques, utilisables plus tard. 

On peut y trouver parfois un peu de traitement de données (i.e sur les résultats).

Il ne faut pas oublier cependant que certaines techniques ou objet technique 
font appel à beaucoup de sciences, et ainsi la description des moyens de mesure employés peuvent
demander des connaissances pointues. 


%*********************************
\subsubsection{Rapport scientifique}
%*********************************

Une entreprise contacte une équipe de rechercher 
pour mener une étude aérodynamique sur le décollement de la couche limite (l'écoulement) autour d'un profil d'aile, en utilisant en partie les moyens
expérimentaux à disposition, comme une soufflerie ou un canal hydrodynamique.

On effectue le montage et les essais.
 En plus de cela est fournie  une bibliographie sur la problématique qui requiert des connaissances scientifiques sur la problématique, et une capacité de synthèse.
  Ensuite  des calculs de couche limite ou des simulations numériques sont effectuées.
  Enfin on doit conduire un post-traitement des données expérimentales en les comparant aux résultats des simulations.
 \'Eventuellement des codes plus ou moins important devront être écrits.   
   Finalement on conclut ce travail en donnant des perspectives sur l'évolution du projet.

On voit ici qu'il y a eu un travail scientifique sérieux qui a été effectué et qui doit être rédigé.
Le rapport sera un rapport scientifique qui contiendra à la fois les aspects expérimentaux, la théorie, les simulations et toutes les analyses produites,
 avec une conclusion sur les aspects scientifiques du problème.

Un rapport scientifique peut se baser sur un rapport technique afin d'avoir des résultats de référence.


%**********************************
\subsection{Rapport de projet}
%*********************************

Dans l'industrie ou la recherche tout travail est lié à un projet
 et donc le rapport de projet est une synthèse  partielle ou complète 
 des actions menées dans le projet, concaténant des informations de rapports techniques et scientifiques éventuellement.
 Il peut être complété par des informations ou des parties traitant des aspects organisationnels et économiques.

\`A l'image de cela, un rapport de projet dans l'enseignement doit contenir un certain nombres d'informations assez similaires à celles d'un projet
scientifique ou industriel.

Une  thèse  de doctorat entre  dans cette catégorie.

Cela sera discuté dans le chapitre \ref{cha:structure}. 


%**********************************
\subsection{Rapport de stage}
%*********************************

Un rapport de stage est formaté comme un rapport de projet. Quelques différences notables sont à ajourer:
\begin{itemize}
\item La page des remerciements est obligatoire, il faut remercier l'entreprise et les collègues qui vous ont aidés
\item Il faut rajouter un chapitre ou une section de 2 à 4 pages qui décrit les activités de l'entreprise et la position du stagiaire au sein de l'entreprise
\item Les attendues du stage doivent être indiquées dans l'introduction
\item Il faut mettre en avant le travail effectué et les apports pour l'entreprise.
\item Il faut éviter, comme dans tous les rapports, un plagiat en reprenant des pages entières de rapport des stagiaires
 qui ont déjà travaillé sur une partie projet. Cela se voit facilement en général et donne une mauvaise impression au lecteur évaluateur.
\end{itemize}

% ************************************
\subsection{Evaluation d'un rapport}
% ************************************

Un document écrit est toujours évalué sur deux aspects:

\begin{itemize}
\item \textbf{la forme} qui correspond à des choix éditoriaux, de style, de mise en forme du texte, des images, des figures et des sections. 
Utiliser des items par exemple est un choix de style qui permet de mettre en évidence des points plus facilement.
\item \textbf{le fond} 
 qui  correspond au texte écrit, ou choix du contenu des figures,
  des tableaux et  des équations présentées.
  
  Le fond couvre tout ce qui dépend  du sujet du rapport.
\end{itemize}


On pourrait  penser a priori que les deux aspects sont totalement découplés et donc en négliger un par rapport à l'autre.
Dans la réalité, il est très rare d'avoir un bon fond quand la forme est très mauvaise. Quelqu'un qui a passé 
du temps à produire un beau document a aussi 
passé du temps à produire un travail scientifique de qualité.

Dans la suite de ce document on parlera uniquement de la forme. 
Le fond de ce document est de donner des lignes de conduite, des guides
afin d'écrire un rapport de qualité sur la forme.


%*****************************************
\section{Outils utilisés pour ce rapport}
%****************************************

Ce rapport a été écrit avec Latex, en choisissant la langue anglaise (ce qui explique que les noms des grandes structures soient en anglais).
Les figures sont souvent générées ou retravaillées avec Inkscape et les courbes sont obtenues en Python.
La bibliographie fait appel à Bibtex et peut être modifiée par exemple avec Jabref. Tous ces outils sont gratuits et sur toutes les plateformes
logiciels habituelles.


%**********************************
\section{Contenu de l'introduction}
%**********************************

\Important{A mettre plus loin ?}

Une bonne introduction est fondamentalement importante pour donner envie au lecteur de lire la suite. C'est le premier regard sur le document, il est 
donc déterminant. 

Il est courant de dire que l'introduction doit être écrite en dernier. Elle doit être relativement courte (au plus 4 ou 5 pages)

Elle  permet d'amener le contexte général et la motivation du projet
notamment en indiquant des applications possibles.

Ensuite elle propose les  objectifs à atteindre 
et  la démarche et les méthodologies employées pour les atteindre.

Enfin  elle expose les résultats attendus.

Des éléments bibliographiques doivent étayer les informations fournies.

L'introduction finalement se conclut en présentant la structure
du document, qui naturellement doit suivre globalement la table des matières. 

Contrairement à ce qui est présenté ici, une introduction peut ne pas contenir de section, mais uniquement un texte avec des paragraphes.

Voici un exemple d'une fin d'introduction qui concerne ce document par ailleurs:

\vs{1.5}

\textit{Le document décrit d'abord dans le chapitre 2 la structure d'un rapport, avec 
les éléments nécessaires et ceux optionnels qui sont présents en fonction du contexte. 
Le chapitre 3 aborde le format des éléments importants d'un document et 
quelques points sur le style. 
Le chapitre 4 est plus focalisé sur le contenu propre à chaque chapitre. 
Des conclusions et perspectives sont données dans le chapitre suivant. Un exemple de bibliographie suit le chapitre 5. 
Le document se termine par un exemple d'annexe.
}